\begin{proposition}[E.2.15; Proved; an isotrivial volume obstruction]
Let $f:X\to\PP^1$ be a proper holomorphic map with connected
fibres from a compact connected complex threefold.  Suppose
that for some $p\in\PP^1$ the restriction over
$U=\PP^1\setminus\{p\}$ is a smooth family of mutually
isomorphic complex two-tori.  If
\[
 K_X^{\otimes A}\simeq f^*\mathcal O_{\PP^1}(D),
 \qquad A\in\Z_{>0},\quad D\in\Z,
 \tag{E.2.15.1}
\]
then necessarily $D<-A$.  In particular the ordinary-source
collapsed-target case of E.2.12 cannot occur for the
six-sphere families under consideration.
\end{proposition}
\begin{proof}
Choose an affine coordinate $t$ on $U\simeq\C$ and put
$q=t^{-1}$ near $p$.  Since $U$ is simply connected, its
integral first-homology local system has a global marking.
The holomorphic period subbundle is constant: locally the
marked periods belong to the countable integral-change-of-basis
orbit of one fixed torus, and a continuous map from a connected
base into such a countable subset of a complex Grassmannian
is constant.  For the last assertion, any two distinct image
points can be separated by a real continuous coordinate;
its connected image would contain an interval and would be
uncountable.  Thus the Jacobian group family is the product
with one fixed torus $T$.

Choose a nonzero translation-invariant holomorphic two-form
$\eta_T$ on $T$.  Local holomorphic origins identify our
family with its Jacobian, and two choices differ by a
translation.  Hence $\eta_T$ defines a global nonvanishing
relative canonical form $\eta$ on $X_U$.  Its fibre volume
is a fixed positive number:
\[
 \int_{X_t}\frac14\eta\wedge\overline\eta=V>0.
 \tag{E.2.15.2}
\]
This construction uses neither a global torsor section nor
a polarization on $X$.  The canonical determinant
isomorphism for a submersion gives the nonvanishing
absolute three-form $\eta\wedge dt$ on $X_U$; this notation
does not require choosing an absolute lift of $\eta$.

Use (E.2.15.1) to pull back a meromorphic section of
$\mathcal O(D)$ whose divisor is $D[p]$.  We obtain a
meromorphic pluricanonical form $\sigma$ with
\[
 \operatorname{div}(\sigma)=D f^*(p).
\]
On $X_U$ it is holomorphic and nowhere zero.  Its quotient
by $(\eta\wedge dt)^{\otimes A}$ is a holomorphic function
constant on each compact connected fibre.  It therefore
descends to an entire nowhere-zero function $\phi$ on $\C$:
\[
 \sigma=\phi(t)(\eta\wedge dt)^{\otimes A}.
 \tag{E.2.15.3}
\]

On a small disc about $p$, the meromorphic section
$q^{-D}\sigma$ is in fact a holomorphic nowhere-zero section
of $K_X^{\otimes A}$ over the whole inverse image.  This
follows directly from (E.2.15.1), or by cancellation of the
actual divisors, with their multiplicities.  Denote it by
$\Omega_A$.  On the punctured inverse image, since
$dt=-q^{-2}dq$, it is
\[
 \Omega_A=(-1)^A P(q)(\eta\wedge dq)^{\otimes A},
 \qquad P(q)=q^{-D-2A}\phi(1/q).
 \tag{E.2.15.4}
\]
The function $P$ is single-valued and holomorphic on the
punctured disc.  No extension of a comparison with another
threefold is asserted here.

A nowhere-zero holomorphic $A$-canonical form determines
a smooth positive density: if it is
$h(z)(dz_1\wedge dz_2\wedge dz_3)^{\otimes A}$, take
$|h|^{2/A}$ times the ordinary complex-coordinate volume
density.  This is independent of the coordinates.  The
density associated with $\Omega_A$ has finite integral
over the inverse image of a smaller closed disc, because
that inverse image is compact.  Fibre integration of
(E.2.15.4), using (E.2.15.2), consequently gives
\[
 \int_{0<|q|<\epsilon}|P(q)|^{2/A}\,d\mathrm{Area}(q)
 <\infty.
 \tag{E.2.15.5}
\]
Only a fixed positive normalization constant is suppressed.
This step concerns a smooth positive density on the actual
total space and requires no K\"ahler metric.

We spell out the isolated-singularity consequence of
(E.2.15.5), including when $2/A<1$.  The function
$|P|^{2/A}$ is subharmonic.  If $|q|=r$, the mean-value
inequality on the disc of radius $r/2$ centred at $q$ gives
\[
 |P(q)|^{2/A}
 \leq \frac{4}{\pi r^2}
       \int_{r/2<|z|<3r/2}|P(z)|^{2/A}\,d\mathrm{Area}(z).
\]
The last integral tends to zero as $r\to0$, by
integrability.  Hence $r^A|P(q)|\to0$, uniformly on
$|q|=r$.  The function $q^A P(q)$ therefore has a removable
singularity and vanishes at zero.  For some holomorphic
$H$ on the disc,
\[
 P(q)=q^{1-A}H(q).
 \tag{E.2.15.6}
\]
Equivalently, the pole order of $P$ is at most $A-1$.

Combining (E.2.15.4) and (E.2.15.6) yields
\[
 \phi(1/q)=q^{D+A+1}H(q).
 \tag{E.2.15.7}
\]
If $D\geq-A$, the right-hand side extends to zero at
$q=0$.  Thus $\phi$ extends holomorphically to $\PP^1$
and vanishes at $p$.  A holomorphic function on $\PP^1$
is constant, so $\phi=0$, contradicting its construction.
This proves $D<-A$.

In E.2.10 the actual canonical-root equality has
\[
 A=12kl,\qquad D=kl-3k-4l,
 \qquad
 \frac DA=\frac1{12}-\frac1{4l}-\frac1{3k}
       \geq-\frac12>-1.
\]
If the collapsed constant source point is ordinary,
E.2.12 says that the central map has precisely one
critical value and is isotrivial on its complement.
All hypotheses above apply to the actual central
threefold, and give the required contradiction.
\end{proof}

\begin{remark}[Scope of the volume argument]
The proof uses the equality of line bundles in
(E.2.15.1), not just the topological vanishing of
$c_1(K_X)$.  It also uses constancy of the torus periods,
which makes the fibre volume in (E.2.15.2) constant.
It does not assert a general good-reduction or
comparison theorem for non-K\"ahler torus degenerations.
Together with the separate elliptic-action exclusion
in E.2.14, it removes the two types of collapsed
target permitted by E.2.12.  Identification of a
noncollapsed interior completion, and the behaviour
on arbitrary new deformation branches, remain
distinct questions.
\end{remark}
